%% T-RO submission — Causal-Validity Auditing for Learned Grasp-Candidate Scoring
%% Compiled with: pdflatex paper_tro.tex (or latexmk -pdf paper_tro.tex)
%% Converted from paper_tro_draft.md -- see that file for the working/annotated version.
%%
%% Required packages (all in standard TeX Live):
%%   IEEEtran, graphicx, booktabs, amsmath, amssymb,
%%   microtype, hyperref, cleveref, xcolor, multirow, array

\documentclass[journal]{IEEEtran}

%% --- Core packages ---
\usepackage[T1]{fontenc}
\usepackage[utf8]{inputenc}
\usepackage{lmodern}
\usepackage{microtype}

%% --- Math ---
\usepackage{amsmath}
\usepackage{amssymb}

%% --- Graphics ---
\usepackage{graphicx}
\graphicspath{{figures/}}

%% --- Tables ---
\usepackage{booktabs}
\usepackage{multirow}
\usepackage{array}

%% --- Color ---
\usepackage{xcolor}
\definecolor{posblue}{HTML}{2166ac}
\definecolor{negred} {HTML}{d6604d}
\definecolor{neutgray}{HTML}{888888}

%% --- Hyperlinks ---
\usepackage[hidelinks]{hyperref}
\usepackage{cleveref}

%% --- Misc ---
\usepackage{balance}

% ─────────────────────────────────────────────────────────────────────────────
\title{Causal-Validity Auditing for Learned Grasp-Candidate Scoring,\\
with a Case Study in Cross-Embodiment Grasping}

% Double-blind: author block omitted, matching paper_final.tex's convention
% \author{...}

% ─────────────────────────────────────────────────────────────────────────────
\begin{document}

\maketitle
\thispagestyle{empty}
\pagestyle{empty}

% =============================================================================
\begin{abstract}
A learned scorer that reranks grasp candidates before execution is valid
only if every input feature is computable before the chosen candidate is
executed -- a constraint we call causal validity, easy to violate silently
since an offline evaluation looks identical either way. We formalize a
pre-execution admissibility criterion, build a reference audit tool, and
show its necessity on our own cross-embodiment grasp-reranking pipeline
(SO-ARM101 to Piper, simulation): a pooling effect replicating five times
at $p{<}0.0001$ collapses to an exact null ($p{=}1.0000$) once invalid
features are removed. An automated static-dataflow tagger independently
catches a contamination bug our own hand-built audit had missed, with a
measurable downstream consequence. Applied prospectively, a predecessor
critic's reported $+15.6$-point gain is shown invalid on two grounds, and
a rebuilt, causally-admissible critic beats a geometric baseline
live-executed on two disjoint batches ($+15.6$/$+14.0$ points,
$p{<}0.003$), replicating on a fourth object added after the design was
frozen while showing a clean zero-shot negative on a fifth. We report a
further independently discovered controller bug as a case study in the
same discipline, eleven ruled-out mechanisms for a still-open reliability
problem, and a hardware deployment architecture not yet physically
validated.
\end{abstract}

\begin{IEEEkeywords}
Grasping, data leakage, reproducibility, static program analysis, robot manipulation, cross-embodiment transfer, benchmarking
\end{IEEEkeywords}

% =============================================================================
\section{Introduction}

Grasp-candidate reranking -- training a lightweight scorer to select the
most promising pose from a pool of candidates before executing any of
them -- is an established pattern in robotic manipulation
\cite{mahler2017dexnet,graspgen2025}. Its correctness rests on an
assumption rarely stated explicitly: every feature the scorer consumes
must be computable \emph{before} the chosen candidate is executed. A
feature that depends on execution's outcome -- contact force during
approach, positional drift during descent, a kinematic residual re-solved
after the arm has already moved -- cannot be known at the moment a
not-yet-executed candidate must be scored, however predictive it looks
offline.

This constraint is easy to satisfy by construction and easy to violate by
accident. GraspGen \cite{graspgen2025} satisfies it by construction: its
discriminator takes only a point-cloud encoding and a candidate pose at
inference time, while training labels come from simulated physics
outcomes. It is violated by accident whenever a feature logged for
training-data purposes gets reused, unexamined, as a live-selection input
-- and the violation is invisible under the evaluation method such systems
reach for first: offline accuracy on historical logs, which already
contain every candidate's outcome and so cannot distinguish a valid
feature from an invalid one.

We hit this failure mode directly while extending a validated
single-embodiment grasp reranker (LGGSN \cite{lggsn_ral_under_review},
pairwise-BPR-trained on a 5-DoF SO-ARM101 arm) to a second, independently
implemented 6-DoF platform (a Piper arm, MuJoCo/RoboSuite simulation). A
cross-embodiment pooling effect replicated across five independent
pilots at $p{<}0.0001$, effect size growing from $+0.21$ to $+0.35$ AUC as
feature sets were enriched. Designing the live-execution test this result
was building toward required, for the first time, identifying which
features would actually be available to a not-yet-executed candidate.
None of the features driving the five-times-replicated result qualified.
Restricting to the causally valid subset collapsed the effect to an exact
null ($p{=}1.0000$).

This paper is organized around that discovery, not around a single
proposed algorithm for improving grasp success. Contributions:

\begin{enumerate}
\item A formal \textbf{pre-execution admissibility criterion}
  (\cref{sec:method}), including a subtlety required in practice: a
  feature may depend on sibling, not-yet-executed pool candidates without
  being causally invalid -- conflating ``depends on multiple candidates''
  with ``depends on execution'' is itself a documented failure mode of our
  own reference implementation.
\item A \textbf{reference audit tool}, enforced as an automatic
  training-time gate, externally validated against two recent published
  systems (\cref{sec:external,sec:retro}).
\item An \textbf{automated dataflow analyzer} inferring feature provenance
  from source code rather than manual annotation, which independently
  caught a contamination bug our manual audit missed and had reported as
  clean -- with a measurable downstream consequence -- and an
  \textbf{interprocedural extension} that resolves a function parameter's
  own provenance by tracing its real call sites, validated on a second,
  independently implemented codebase (\cref{sec:autotag}).
\item A \textbf{second, independent case study} in the same discipline: a
  silent double-scaling bug in a widely used composite robot controller,
  validated on three previously near-zero-success object categories
  (\cref{sec:gripperbug}).
\item An \textbf{honest, statistically rigorous ledger of eleven
  ruled-out mechanisms} for the execution-time reliability problem our
  diagnostics converged on, none presented as solved
  (\cref{sec:negative}).
\item A \textbf{hardware deployment architecture and safety checklist},
  scoped plainly as not yet physically validated (\cref{sec:hardware}).
\end{enumerate}

We do not claim to have solved cross-embodiment grasp transfer or the
execution-precision problem our diagnostics repeatedly point at. We claim
that trying to solve them surfaced a general, underappreciated evaluation
failure mode with a formal characterization and a checkable fix -- and
that even a careful team applying that fix to its own work will still
find real instances of it, which is worth reporting rather than hiding.

% =============================================================================
\section{Related Work}
\label{sec:related}

\textbf{Grasp candidate scoring.} Dex-Net 2.0 \cite{mahler2017dexnet}
trains a grasp-quality CNN offline on millions of simulated grasp
outcomes, then deploys the trained network -- not a live simulation query
-- to rank real candidates from a depth image at inference time;
execution-derived information appears only as training supervision.
GraspGen \cite{graspgen2025} follows the same pattern with an
on-generator-trained discriminator whose inference-time inputs are a
point-cloud encoding and a candidate SE(3) pose. Neither work states the
input/label distinction as a formal, checkable criterion; both satisfy it,
which we use as external validation of our formalization
(\cref{sec:external}) rather than as prior art for the criterion itself.

\textbf{Path-aware grasp feasibility.} Recent work on grasp-dependent
feasibility prediction \cite{feasibility2025} generates causally
admissible, path-level geometric labels -- inverse-kinematics feasibility
and a transit-collision flag from mesh sweeps along a fixed waypoint
template -- and trains a compact MLP on them for candidate ranking,
validated as transferring to physics-executed trajectories on a single
rigid cuboid with a fixed waypoint template. This is framed entirely as an
efficiency argument (cheap labels vs.\ expensive physics simulation), with
no connection to a causal-validity theory. We considered adapting this
mechanism for our own still-open execution-reliability problem and
rejected it (\cref{sec:negative}) on the grounds that our diagnosed
failure mode (sustained, necessary contact under sub-millimeter
miscentering) does not match the binary avoidable-collision framing this
mechanism targets, and is mechanistically closer to four contact-response
interventions we had already tested and found ineffective.

\textbf{Cross-embodiment transfer in grasping.} GraspGen-X
\cite{graspgenx2026} demonstrates cross-embodiment 6-DoF diffusion-based
grasping validated on a real AgileX Piper arm, using pointwise (not
pairwise) supervision. Freeform Preference Learning \cite{fpl2026} learns
a language-conditioned reward model from trajectory-level human
preferences for post-hoc reward shaping in reinforcement learning -- a
role that legitimately requires execution-derived input by design, which
we use as a boundary case our criterion correctly does not flag
(\cref{sec:external}).

\textbf{Data leakage and offline evaluation.} Look-ahead and label
leakage are recognized general problems in machine learning and have been
studied recently in offline reinforcement-learning policy evaluation
\cite{offlinerl_leakage2026}, where an oracle policy may exploit
information unavailable to the deployed agent. Closest to our formal
criterion is concurrent work by Roth \cite{rothgrammar2026}, a grammar of
eight typed primitives with a call-time-enforced admissibility gate that
makes train/test leakage structurally unrepresentable in standard tabular
and time-series ML pipelines. Two differences are worth stating plainly.
First, domain: Roth's boundary is the train/test split; ours is whether a
candidate has physically executed yet, including a pool-relative case
(a feature may depend on sibling, not-yet-executed candidates) with no
analogue in a train/test framework. Second, mechanism: Roth's grammar
requires authoring a pipeline against his primitives (reference
implementations for Python and R); our tool statically analyzes arbitrary,
already-written code via a single per-function marker, requiring no
pipeline rewrite. Also relevant, at the level of general practice rather
than a formal criterion: LeakageDetector \cite{leakagedetector2025} is a
PyCharm plugin that flags generic train/test-split leakage patterns in
tabular ML code via static analysis, and closely related notebook-specific
work does the same for Jupyter pipelines -- both operate on the same
generic train/test boundary as Roth, not the temporal,
already-executed-or-not boundary this domain requires, and neither targets
robot-manipulation pipelines specifically. Concurrent work by Jiang et
al.~\cite{benchmarkaudit2026} audits five robot-manipulation benchmarks
(LIBERO, CALVIN, SimplerEnv, RoboCasa, RoboTwin~2.0) for four largely
orthogonal failure modes -- shortcut solvability, insufficient statistical
power, creeping overfitting, and data-source dependence -- via empirical
probes (a weak language-blind model matching reported SOTA on LIBERO;
performance collapse under in-distribution pose randomization on CALVIN),
not static code analysis, and does not check whether individual training
or scoring features are computable before an action is taken. We view this
as a complementary, not competing, direction: their diagnostics answer
whether a benchmark's *reported numbers* are trustworthy; ours answers
whether a *specific pipeline's input features* could not have been known
before the choice they inform, a narrower and, in our experience
(\cref{sec:external}), more direct route to the specific failure mode this
paper's central finding depends on. We are not aware of prior work stating
our specific criterion for grasp-candidate scoring features, nor of an
automated tool for detecting the violation from unmodified source code in
this domain.

\textbf{Perception uncertainty for grasping.} Closest to our planned
real-noise-characterization extension (\cref{sec:conclusion}) is
Consensus-Driven Uncertainty \cite{consensusuncertainty2025}, which trains
an MLP to predict grasp failure from disagreement among three independent
pretrained pose estimators (EPOS, GDRNPP, ZebraPose) applied to the same
real image, and reports that "poses perturbed by Gaussian noise will not
fail the same way that real pose estimators fail" -- independent evidence
for the same hypothesis motivating our own noise-characterization plan.
Two gaps remain open relative to our design: that work does not
quantitatively characterize the noise's statistical structure (per-axis
variance, Gaussianity, position-orientation correlation), and its
uncertainty signal drives an accept/reject gate on a single proposed
grasp rather than a selection rule over a candidate pool.

\textbf{Contact-aware execution.} CoorGrasp \cite{coorgrasp2026} and
related work on wrench-balance-criterion phase transitions use real
tactile sensing to coordinate arm motion after fingers have already
formed contact with an object. Lacking tactile hardware, we used
MuJoCo's own contact-force API as a privileged, hardware-free proxy for
this class of mechanism (\cref{sec:negative}), finding it -- along with
three related contact-response variants -- ineffective for our specific
failure mode.

% =============================================================================
\section{System}

Our platform is a 6-DoF AgileX Piper arm simulated in MuJoCo via
RoboSuite, a second, independently implemented embodiment alongside this
project's original 5-DoF SO-ARM101 platform. The core execution primitive
drives the arm through a fixed phase sequence -- transit, approach,
descend, close, lift, transit-to-tray, lower, retract -- using damped
least-squares inverse kinematics solved against a saved kinematic
snapshot, and an absolute joint-position controller. Success is defined as
final object position within a fixed radius of the tray center.

LGGSN, this project's validated single-embodiment contribution, is a
lightweight MLP-based pairwise reranker trained with Bayesian Personalized
Ranking loss over a 14-dimensional geometric feature space on the SO-ARM101
platform, and is the deployed, live candidate scorer for that platform's
production pipeline. It is both reused infrastructure in this paper and,
in \cref{sec:method}, the system under audit -- the causal-validity
criterion is applied directly to its live, already-deployed feature
pipeline, not to a toy example.

% =============================================================================
\section{Causal-Validity Auditing for Learned Grasp-Candidate Scorers}
\label{sec:method}

\subsection{Definitions}

\textbf{Definition 1 (Candidate pool, execution trace).} A scene induces a
candidate pool $C = (c_1,\ldots,c_N)$, each candidate specified by a pose
$p_i \in SE(3)$ and generator-time attributes $a_i$. Let $G$ denote
static, execution-independent scene/object information. If $c_i$ is
physically executed, let $T_i$ denote its execution trace -- contact
forces, intermediate poses, net displacement, and the terminal success
label $y_i \in \{0,1\}$.

\textbf{Definition 2 (Selection-time information set).} At the moment a
scorer must choose which candidate(s) to execute, the selection-time
information set is
\[
I(C) = \{p_1,\ldots,p_N\} \cup \{a_1,\ldots,a_N\} \cup G \cup K(C,G),
\]
where $K(C,G)$ is the set of values obtainable from isolated,
execution-free queries against a candidate pose and $G$ (e.g.\ a
kinematic IK residual or collision check evaluated without moving the
arm). $I(C)$ contains no $T_j$ for any $j$ -- nothing in the pool has
executed yet.

\textbf{Definition 3 (PRE\_EXECUTION-admissible feature).} A feature
function $\varphi$ is admissible if there exists $g$ such that, for every
candidate $c_i$ in every pool $C$, $\varphi(c_i) = g(I(C))$. $I(C)$
includes every candidate's pose, not just $c_i$'s own -- a feature may
depend on sibling candidates (e.g.\ distance to the pool centroid) without
being execution-derived. Conflating ``depends on more than one
candidate'' with ``depends on execution'' is a documented failure mode of
our own reference implementation (\cref{sec:selfcorrections}). Otherwise
$\varphi$ is EXECUTION\_DERIVED.

\textbf{Proposition.} A scorer $S(c_i) = h(\varphi_1(c_i),\ldots,\varphi_k(c_i))$
used to select among not-yet-executed candidates is causally valid iff
every $\varphi_j$ is admissible. A feature set containing an
execution-derived $\varphi_j$ is invalid regardless of how well $S$
performs offline on historical $(c_i, T_i)$ pairs -- offline evaluation
implicitly supplies $T_i$ for every evaluated $i$, satisfying $\varphi_j$
even though it would be undefined for a genuinely new candidate.

\textbf{Corollary.} Execution-derived quantities remain valid as
supervision targets or as values used only to construct them -- the
constraint binds the input side of $S$, not the label side. This is why
GraspGen's design (\cref{sec:related}) is correct.

\textbf{Algorithm.} (1) Tag every raw field with provenance
$\tau(f) \in \{\mathrm{PRE}, \mathrm{EXEC}\}$ by tracing the code path
that computes it, not a comment or an assumption about a similarly named
field. (2) Before training or invoking any live-selection scorer with
feature set $F$, reject if any $f \in F$ has $\tau(f)=\mathrm{EXEC}$ or is
unregistered -- fail closed. (3) Re-run (2) whenever $F$ changes or
whenever the code computing any $f \in F$ changes, since $\tau(f)$ is a
property of code, not of a field's name.

\subsection{External grounding}
\label{sec:external}

We checked the criterion against two recent published systems chosen to
stress-test both directions of the boundary. GraspGen \cite{graspgen2025}
is confirmed to satisfy it: its discriminator's inference-time inputs are
a point-cloud encoding and candidate pose only; training labels are
simulated shaking-motion outcomes. Freeform Preference Learning
\cite{fpl2026} is confirmed as a correct \emph{non-flagged} boundary case:
its reward model legitimately takes trajectory-level, execution-derived
input, because its role in the pipeline -- post-hoc reward shaping for
reinforcement learning -- is categorically different from live
pre-execution candidate selection. The criterion correctly distinguishes
``execution-derived input because the model's role requires it'' from
``execution-derived input leaked into a model whose role requires
admissibility,'' which is the failure mode we hit.

\subsection{Retrospective demonstration}
\label{sec:retro}

We implemented the registry (field-level provenance tags with an
enforced audit function) covering every raw field logged by both
platforms' pipelines and ran it against every feature set used across
this project's cross-embodiment reranking pilots, in chronological order.
The audit correctly flags every pilot that produced the illusory
pooling-beats-zero-shot effect and passes every pilot ultimately trusted;
it attributes the contamination specifically to the Piper side (a
descend-phase IK residual logged after physical motion, and a
positional-drift measurement by definition only available post-motion),
while the SO-ARM101 side's superficially similar fields -- traced against
the actual live inference code rather than a comment describing an
inactive legacy dataset -- turned out to be a harmless pre-execution proxy
plus dead constant-zero features.

\subsection{Self-corrections}
\label{sec:selfcorrections}

Building this registry produced two errors of its own, both caught by
re-tracing against live code rather than trusting an initial plausible
explanation: (1) the SO-ARM101 fields above were initially mislabeled
execution-derived by quoting a comment describing a different, unused
dataset; (2) two pool-relative features were initially mislabeled by
conflating ``depends on sibling not-yet-executed candidates'' with
``depends on execution.'' We report both because they demonstrate the
discipline the method argues for, applied to itself, not because they
weaken the tool's credibility.

\subsection{Automated tagging}
\label{sec:autotag}

A manually populated registry does not scale and, as
\cref{sec:selfcorrections} shows, is itself error-prone. We built a
static dataflow analyzer that infers provenance automatically: a
fixed-point pass over the module identifies every \emph{execution-touching}
function (one that directly or transitively calls the simulator's
physical-actuation entry point), and a forward taint-propagation pass over
the target function's body -- gated by a single human-placed marker
denoting the point after which the current candidate has begun physically
executing -- determines, for every field in the function's output,
whether it is reachable from any post-marker execution-touching call or
live-state read. This reduces the manual annotation burden from
$O(\text{number of logged fields})$ to $O(\text{number of physical commit
points in the codebase})$ -- a single marker, in our case, rather than
thirteen-plus per-field judgments.

Run against the real, live pick-and-place function, the tagger
independently disagreed with our hand-built registry on one field: it
flagged a grasp-orientation feature as execution-derived, where the
manual registry had marked it admissible. Tracing why confirmed the
tagger was correct: the underlying pose variable is reassigned twice --
once pre-commit during candidate selection (admissible) and once
post-commit, at a deliberate mid-trial drift-correction step that
re-reads simulator state after the descend phase has already executed. A
human reading the function once sees the first assignment and does not
track the second; the automated tagger mechanically follows every
reassignment. This was not a cosmetic error: the field was part of the
feature set our own prior retrospective validation
(\cref{sec:retro}, ``corrected'' configuration) had already reported as
clean. Re-running that validation with a genuinely admissible two-feature
set (dropping the contaminated field from both platforms for a fair
shared feature space) confirmed the qualitative finding is unchanged --
all pooling conditions remain exactly identical to zero-shot transfer
($\mathrm{diff}=0.0000$) -- but the reported absolute accuracy changed
from $0.8236$ to $0.1327$, a number that had already been published as
final in this project's own working documentation before the automated
check caught it.

We verified, rather than assumed, why: the object top-surface offset
feature is an exact constant across every row of the single-object
(Cracker-only) Piper collection ($\sigma{=}0$ over $n{=}250$), and the
remaining pose-height feature is a near-constant \emph{per scene} --
$14/25$ scenes show within-scene spread under $2{\times}10^{-4}$,
consistent with floating-point/physics-settling noise rather than a
genuine per-candidate signal, because it is read once, from the object's
own spawn pose, before any candidate-specific action -- identical
regardless of which of the ten pooled candidates is later attempted.
Together this means the corrected two-feature set has almost no capacity
to discriminate between different candidates drawn from the \emph{same}
scene, which is exactly the comparison the pairwise ranking objective is
trained and evaluated on. This explains the exact equality across all
four pooling conditions (there is essentially nothing in the input for
any of them to learn from) and reframes $0.1327$ not as an unexplained
oddity but as an expected consequence of a feature set that is causally
valid yet structurally near-uninformative for this specific comparison --
a genuinely predictive but still-admissible per-candidate feature (e.g.\
the original, pre-commit candidate orientation, distinct from the
execution-derived value \cref{sec:autotag} identified) was not tested
here and is left as future work.

We validated the tagger against a second, structurally different function
(a pure data-serialization routine with no internal execution boundary
and a constructor-call field pattern rather than a return literal) and
initially found a genuine scope limitation worth stating rather than
hiding: because this function contains no physical-commit marker of its
own, every field trivially returns admissible, which happens to be
correct (verified independently against the function's known input
provenance) but was not actually \emph{demonstrated} by the tool -- the
analysis was single-function-scoped and could not verify a parameter's
own provenance across a call boundary.

We closed this by extending the tool with interprocedural resolution: for
a target parameter, find every real call site in the module and evaluate
the taint of the bound argument at that call site, using the calling
function's own taint state up to that point, rather than assuming the
parameter is safe. Two fixes were necessary on this second, independently
implemented (PyBullet-based, not MuJoCo-based) codebase: the
physical-actuation entry-point check, previously hardcoded to a single
method name, had to be generalized to a configurable set, since this
codebase's execution calls are named differently; and the calling
function itself needed its own commit marker, since -- consistent with
the design already established -- without one the tagger correctly flags
nothing as tainted rather than guessing. This generalizes the earlier
finding: interprocedural resolution does not shrink the marker
requirement to one marker per codebase, it requires one marker per
function with a genuine execution boundary, applied transitively up the
call chain. With both fixes in place, resolving the real call site
correctly and automatically recovered every field's true provenance --
including confirming that a success-derived label field is
execution-derived (correct, since it is the training label, exempt under
the Corollary) while the pose-derived features are genuinely
pre-execution and safe for live selection.

\subsection{An independently discovered infrastructure bug, and a stronger re-verification}
\label{sec:rngbug}

While instrumenting a new admissible feature to test directly (whether the
original, pre-commit candidate orientation -- distinct from the
disqualified execution-derived value in \cref{sec:autotag} -- carries
real signal), we found a second, unrelated bug: the Piper simulation's
object-placement sampler was constructed without an explicit random-number
generator, silently falling back to an OS-entropy-seeded default
completely independent of the \texttt{numpy} legacy seed this project's
every prior experiment relied on for reproducibility. Concretely, a
scene's ten pooled candidates were never actually facing the same object
placement, despite the code's own documentation asserting they were;
confirmed empirically by re-running an unmodified data-collection script
on an identical scene twice and observing three different outcome
patterns across three runs, and independently by finding 7--9\,cm of
spawn-position spread within already-collected data that should have been
constant. We fixed this by explicitly deriving a seeded generator from the
legacy random state at scene-construction time, verified by confirming
bit-identical object placement across repeated runs of the same seed.

This bug does not affect \cref{sec:external,sec:retro,sec:selfcorrections,sec:autotag}: the causal-validity criterion and every
provenance verdict in this paper come from tracing code, not from
statistical properties of collected data. It does affect any specific
accuracy number computed on data grouped by an assumption of shared
placement. We therefore re-collected the 25-scene Piper dataset under the
fix and re-ran the corrected two-feature evaluation: the qualitative
finding is unchanged -- zero-shot, pooled-\texttt{none},
pooled-\texttt{additive}, and pooled-\texttt{interaction} remain exactly
identical ($0.1036$, $\mathrm{diff}=0.0000$) -- confirming the paper's
core claim is not an artifact of the placement bug.

We then tested the untested admissible feature this section had flagged
as future work. A naive pooled check looked dramatic: candidate
orientation correlates with success at $r=-0.55$ ($p{<}0.0001$, $n{=}250$)
-- far stronger than any other signal in this investigation. Decomposing
it, following this paper's own stated discipline rather than reporting
the encouraging number, revealed why: mean orientation and success rate
correlate strongly \emph{between} scenes ($r=-0.71$, $p{=}0.0001$ across
25 scenes) because different object placements are both associated with
systematically different orientations and, for unrelated geometric
reasons, different difficulty -- but the \emph{within}-scene correlation,
the only one that could actually inform a live selection decision among
candidates facing the same already-placed object, is indistinguishable
from zero (mean $+0.03$ across 10 mixed-label scenes, no individual scene
reaching significance). A feature that appears, by a wide margin, to be
the strongest predictor found in this entire investigation carries zero
real selection-relevant signal once correctly decomposed -- a sharper
demonstration of this paper's thesis than the original finding it was
built to test.

\subsection{Prospective application: an object-relative counterfactual grasp critic}
\label{sec:critic}

\Cref{sec:retro,sec:selfcorrections,sec:rngbug} show the criterion catching
mistakes \emph{after} they were made. This section applies it \emph{from
the start}, on a different learned-critic pipeline (\texttt{world\_model/},
predating \cref{sec:method} entirely) that reported a $+15.6$ percentage
point improvement over a geometric heuristic. Reproducing this claim
directly against the original data, we found it invalid on two grounds,
neither of which is itself a causal-validity violation. First, the
evaluation harness's random seed encoded the compared method's identity, so
paired ``geometry'' and ``world-model'' trials never shared a scene or
candidate pool: of 250 nominally paired trials, zero shared even the
executed candidate position. Second, the success criterion counted
transient post-close gripper contact, checked before any lift, as success
--- a criterion that saturates to $150/150$ regardless of candidate
quality, independent of the pairing defect. Correcting the harness alone,
this study's execution-density (roughly thirteen executions per scene, an
order of magnitude above prior usage) surfaced a further,
shared-infrastructure bug: the production
\texttt{physics\_weld\_after\_bilateral} grasp primitive's weld-attach gate
accepted single-jaw contact, not the bilateral contact its name and
protocol require. We fixed it in \texttt{env\_soarm.py}; we flag, but do
not re-verify here, that this may affect other results computed with the
same primitive elsewhere in this project.

We rebuilt the critic to be causally admissible by construction: candidate
pose expressed relative to the target object rather than in world frame,
point-cloud statistics, and object identity --- every field registered
\texttt{PRE\_EXECUTION} in \texttt{provenance.py} and checked by the same
gate as \cref{sec:method}'s other pipelines --- trained with a within-scene
Bradley--Terry/BPR-style pairwise term over successful and failed
candidates from the same scene, plus binary cross-entropy on outcome. Under
the corrected, paired, live-executed harness, the original checkpoint fails
a pre-registered gate outright: AUROC $=0.4996$ on 1{,}500 real
per-candidate labels, chance level, fully explained by contamination from
the criterion in the paragraph above rather than left as a mystery. The
rebuilt critic, evaluated live --- not offline-rescored --- on two scene
batches disjoint from training, from model selection, and from each other
(disjointness independently key-intersection-verified, not asserted),
beats the geometric baseline on both (\cref{tab:critic_results}):
$+15.6$ points on an independent
development-test batch ($44/90$ vs.\ $30/90$; McNemar's exact
$p=0.00258$), and $+14.0$ points on a frozen confirmatory batch never
inspected before this gate was evaluated ($75/150$ vs.\ $54/150$; 27 wins
to 6 losses; McNemar's exact $p=3.24\times10^{-4}$). A small,
physically-real non-determinism footnote applies to both: repeated
identical-pose executions flip their boolean outcome at roughly
$0.6$--$1\%$, the contact solver not being bit-reproducible on marginal
grasps --- small enough not to change either conclusion, but stated rather
than assumed away.

\begin{table}[t]
\centering
\caption{Live-executed critic results (\cref{sec:critic}). The stale
checkpoint is evaluated under the same corrected, paired harness as the
rebuilt critic --- the harness fix alone does not save it.}
\label{tab:critic_results}
\small
\begin{tabular}{@{}lrrrr@{}}
\toprule
Batch & $n$ & Geometry & Critic & McNemar exact $p$ \\
\midrule
Dev-test & 90 & $30/90$ (33.3\%) & $44/90$ (48.9\%) & $0.00258$ \\
Confirmatory (frozen) & 150 & $54/150$ (36.0\%) & $75/150$ (50.0\%) & $3.24{\times}10^{-4}$ \\
\bottomrule
\end{tabular}
\\[0.3em]
{\footnotesize Stale checkpoint (pre-audit), same corrected harness: AUROC $=0.4996$ on 1{,}500 real per-candidate labels --- chance level.}
\end{table}

Consistent with this paper's discipline of reporting negative and
incomplete findings with the same rigor as positive ones, three results
from the same investigation do not support further claims
(\cref{tab:critic_negative}). An
ensemble-uncertainty risk gate, calibrated on held-out data, adds no
measurable benefit over the ungated critic: its coverage is $98.7\%$, i.e.\
it almost never fires. Whether the within-scene pairwise term is
independently responsible for the gain, versus the object-relative feature
representation alone, is unresolved: object-relative BCE with no pairwise
term and the full counterfactual variant are statistically
indistinguishable on the frozen confirmatory batch (2 wins, 3 losses,
$p=1.0$) --- this paper's novelty claim in \cref{sec:critic} rests on
object-centric, causally-admissible scoring beating geometry, not on the
pairwise mechanism specifically. Finally, a small-data (15 demonstrations)
ACT imitation-policy pilot integrates end-to-end with this pipeline but
fails its first closed-loop rollout, knocking the target object off
position; five demonstrations per object is not enough data for any
closed-loop recovery behavior to have been learned. We report this as
integration status, not as a result bearing on \cref{sec:critic}'s central
claim.

\begin{table}[t]
\centering
\caption{Negative and incomplete results from the same investigation (\cref{sec:critic}), reported with equal rigor.}
\label{tab:critic_negative}
\small
\begin{tabular}{@{}p{0.24\linewidth}p{0.68\linewidth}@{}}
\toprule
Item & Result \\
\midrule
Uncertainty risk gate & No measurable benefit; coverage $98.7\%$ (almost never fires) \\
Pairwise-loss contribution & Unresolved: BCE-only vs.\ full counterfactual indistinguishable (2W/3L, $p{=}1.0$) \\
ACT imitation pilot (15 demos) & Fails first closed-loop rollout (single online trial, $dz{=}{-}0.0451$) \\
\bottomrule
\end{tabular}
\end{table}

\subsection{Mechanistic analysis: multi-head contact/lift/success decomposition}
\label{sec:multihead}

\Cref{sec:critic} establishes \emph{that} the corrected critic beats
geometry, live-executed --- this paper's primary quantitative claim. This
section asks \emph{why and where it works}, via a second, complementary,
explicitly \emph{offline re-scoring} experiment that must not be merged
with or presented as confirming \cref{sec:critic}'s live-executed numbers.
Using the same causally-admissible features, we train four prediction
heads instead of one: binary \texttt{bilateral\_contact}, binary
\texttt{lifted}, binary \texttt{success}, and a three-class
\texttt{failure\_type} (\texttt{success} / \texttt{no\_contact} /
\texttt{weld\_no\_lift}). Two further classes originally specified in this
taxonomy, \texttt{contact\_no\_weld} and \texttt{lifted\_then\_dropped},
were confirmed structurally absent from every collected batch: the weld
gate is currently identical to bilateral contact by construction, and the
post-grasp settle window's \texttt{fell\_off} threshold never triggers.
The three-class scheme is what was actually trained; the six-class taxonomy
remains documented future schema, not a current claim. Training uses the
same base-100/dev-200/confirmatory-300 split chain as \cref{sec:critic},
with checkpoint and loss-weighting choice frozen on dev-200 only and
confirmatory-300 read exactly once (audit trail:
\texttt{confirmatory\_run\_log.jsonl} confirms a single run).

On the confirmatory-300 batch, the success head alone, used for top-1
selection, beats geometry pooled by $+10.7$ points ($46.0\%$ vs.\ $35.3\%$;
McNemar's exact $p=0.0052$) --- directionally and in rough magnitude
consistent with \cref{sec:critic}'s live-executed $+14.0$ points, reported
here as separate, complementary evidence for the same claim, not pooled or
averaged with it (\cref{tab:multihead_vs_geom}). All three heads reach AUROC $0.90$--$0.93$ with
$\mathrm{ECE}\approx 0.04$ (\cref{tab:multihead_auroc}). The decomposition's interpretability payoff is
a specific, actionable asymmetry a single-head success critic cannot show:
the model is conservative, not dangerously overconfident, on success --- it
achieves $92.2\%$ recall on true \texttt{no\_contact}, but misclassifies
$32\%$ of true successes as \texttt{no\_contact} (\cref{tab:confusion}). One limitation is
load-bearing, not a footnote: every \texttt{weld\_no\_lift} training
example ($118/118$) is drill-attributed; excluding the drill collapses
that class's support to a single example. This must be read as a
within-drill measurement, not a demonstrated cross-object generalization
capability for that failure mode; a genuine leave-one-object-out test
was not run in this pass. A methodological note we keep rather than hide:
an \texttt{equal}-vs.-\texttt{success\_weighted} loss-weighting ablation
produced measurably different trained parameters, confirmed by direct
comparison, but \emph{identical} argmax top-1 accuracy at every one of 5
seeds on the small internal validation split used during training --- only
the larger dev-200 batch's continuous AUROC distinguished the two settings.
Coarse top-1 metrics can be an insufficiently sensitive selection criterion
for this class of ablation at small validation-set scale; this was caught
by cross-checking against a second metric, not silently reported as ``no
difference.'' Neither this section nor \cref{sec:critic} licenses a claim
about risk-gate or imitation-policy performance --- \cref{sec:critic}'s own
negative results on both stand unchanged; the gains reported here concern
the critic's own outputs, not downstream gating or policy performance.

\begin{table}[t]
\centering
\caption{Success head vs.\ geometry, offline re-scoring, confirmatory-300 (\cref{sec:multihead}). Cracker's exact tie reproduces the same object-level pattern seen in \cref{sec:critic}'s independent live-executed evaluation.}
\label{tab:multihead_vs_geom}
\small
\begin{tabular}{@{}lrrrr@{}}
\toprule
& Geometry & Critic (success head) & $\Delta$ & McNemar exact $p$ \\
\midrule
Pooled ($n{=}150$) & 35.3\% & 46.0\% & $+10.7$pp & $0.0052$ (23W/7L) \\
Cracker ($n{=}50$) & 18.0\% & 18.0\% & $0.0$pp & $1.0$ (exact tie) \\
Mustard ($n{=}50$) & 58.0\% & 80.0\% & $+22.0$pp & $0.0074$ \\
Drill ($n{=}50$) & 30.0\% & 40.0\% & $+10.0$pp & $0.302$ (n.s.) \\
\bottomrule
\end{tabular}
\end{table}

\begin{table}[t]
\centering
\caption{Per-head discrimination and calibration, confirmatory-300 (\cref{sec:multihead}).}
\label{tab:multihead_auroc}
\small
\begin{tabular}{@{}lrrr@{}}
\toprule
Head & AUROC & AUPRC & ECE \\
\midrule
\texttt{bilateral\_contact} & 0.934 & 0.861 & 0.037 \\
\texttt{lifted} & 0.902 & 0.749 & 0.039 \\
\texttt{success} & 0.903 & 0.749 & 0.040 \\
\bottomrule
\end{tabular}
\\[0.3em]
{\footnotesize Per-object AUROC for \texttt{lifted}/\texttt{success}: cracker $0.977$, mustard $0.92$--$0.923$, drill $0.788$ --- an in-distribution difficulty gradient (all three objects were trained on), not an OOD finding.}
\end{table}

\begin{table}[t]
\centering
\caption{\texttt{failure\_type} confusion matrix, pooled confirmatory-300 (\cref{sec:multihead}). Rows are true labels.}
\label{tab:confusion}
\small
\begin{tabular}{@{}lrrr@{}}
\toprule
True $\backslash$ Pred & success & no\_contact & weld\_no\_lift \\
\midrule
success ($n{=}339$) & 212 (62.5\%) & 109 & 18 \\
no\_contact ($n{=}1043$) & 58 & 962 (92.2\%) & 23 \\
weld\_no\_lift ($n{=}118$) & 19 & 10 & 89 (75.4\%) \\
\bottomrule
\end{tabular}
\end{table}

\begin{table}[t]
\centering
\caption{\texttt{weld\_no\_lift} object attribution, train-100 (\cref{sec:multihead}) --- the basis for this section's load-bearing limitation.}
\label{tab:weldnolift}
\small
\begin{tabular}{@{}lrrr@{}}
\toprule
Object & success & no\_contact & weld\_no\_lift \\
\midrule
Cracker & 64 & 335 & 1 \\
Mustard & 108 & 292 & 0 \\
Drill & 120 & 195 & 85 \\
\midrule
All three (used) & 292 & 822 & 86 \\
Excluding drill & 172 & 627 & \textbf{1} \\
\bottomrule
\end{tabular}
\end{table}

\subsection{Generalization to a fourth object, and a clean out-of-distribution boundary}
\label{sec:fourobj}

\Cref{sec:critic,sec:multihead} establish the causally-admissible critic on
three objects (cracker, mustard, drill). This section asks whether the
result generalizes to an object added after the critic's design was frozen,
and whether it generalizes zero-shot to an object never trained on at all
--- two different, and differently answered, questions. We added a fourth
object, a real YCB Tomato Soup Can (internally keyed as
\texttt{cylinder}, a pre-existing asset alias, not a substitute object),
using an independent train/dev/confirmatory seed chain (701/801/901)
disjoint from the three-object study. A YCB banana was considered and
excluded before any critic training: a controlled asset sanity check
produced zero bilateral contact under the current grasp protocol, an
asset-geometry issue orthogonal to this section's claim. A YCB pear was
held out entirely as a zero-shot out-of-distribution (OOD) probe --- never
in the training, dev, or confirmatory pool for any checkpoint below.

Retrained with a four-way object one-hot (three-object checkpoints and
training code untouched), the isolated four-object critic reaches AUROC
$=0.8903$, Brier $=0.1076$, and $\mathrm{ECE}=0.0304$ on an independent
confirmatory set (offline re-scoring, 200 scenes, disjoint from dev
selection --- the same methodology as \cref{sec:multihead}, not
\cref{sec:critic}'s live-executed evaluation), with pooled offline top-1
$52.5\%$ (cracker $18\%$, mustard $78\%$, drill $28\%$, Tomato Soup Can
$86\%$).

A separate, live-executed neighborhood-robustness evaluation (candidate
selection under five local pose/yaw perturbations per scene, 30
scenes/object, 120 scenes total, \cref{tab:fourobj_neighborhood}) beats
geometry pooled for every critic variant (McNemar's exact
$p{<}1{\times}10^{-6}$, Holm-adjusted). The gain is object-dependent:
strong for cracker and drill, significant but smaller for Tomato Soup Can
($p=0.0113$), and not significant for mustard (already near ceiling for
geometry alone at this task). Evaluated zero-shot on the held-out pear ---
an all-zero object one-hot, never seen in training --- every critic
variant \emph{loses} to geometry (point $28.7\%$, mean $27.3\%$,
worst-case $28.0\%$ vs.\ geometry's $40.0\%$; McNemar's exact
$p<0.02$ for all three, favoring geometry). We report this as a clean,
explicit negative boundary, not pooled with the in-distribution result:
this critic improves candidate selection and local-perturbation robustness
within its trained object distribution, and does not presently generalize
zero-shot beyond it.

\begin{table}[t]
\centering
\caption{Live-executed neighborhood-robustness selection, four-object critic vs.\ geometry (\cref{sec:fourobj}), and the zero-shot pear OOD probe (last row, not pooled with the rest).}
\label{tab:fourobj_neighborhood}
\small
\begin{tabular}{@{}lrrrr@{}}
\toprule
Object ($n{=}150$ each) & Geometry & Point & Mean & Worst-case \\
\midrule
Cracker & $7.3\%$ & $29.3\%$ & $28.0\%$ & $30.7\%$ \\
Mustard (n.s., all three) & $83.3\%$ & $80.0\%$ & $80.0\%$ & $83.3\%$ \\
Drill & $43.3\%$ & $66.0\%$ & $64.7\%$ & $66.7\%$ \\
Tomato Soup Can ($p{=}0.0113$) & $69.3\%$ & $80.0\%$ & $80.0\%$ & $80.0\%$ \\
Pooled ($n{=}600$, $p{<}10^{-6}$) & $50.8\%$ & $63.8\%$ & $63.2\%$ & $65.2\%$ \\
\midrule
Pear OOD ($n{=}150$, zero-shot) & $\mathbf{40.0\%}$ & $28.7\%$ & $27.3\%$ & $28.0\%$ \\
\bottomrule
\end{tabular}
\\[0.3em]
{\footnotesize Row-level $p$-values are object-stratified tests over all three critic variants jointly, not per-column; pooled $p$-values (bottom two rows) are Holm-adjusted paired McNemar per variant, all $p{<}10^{-6}$. Bold marks the winning column for pear, the one row where geometry beats every critic variant ($p<0.02$ for all three), reported as an explicit negative boundary.}
\end{table}

A further question is whether a learned auxiliary head that directly
predicts local-perturbation survival (\texttt{local\_perturbation\_survival},
trained on an exhaustive candidate-level label set that executes every
candidate rather than only the selected one, removing the selection bias
this project's own audit criterion exists to catch) adds value beyond the
neighborhood-aggregation methods already reported above. Offline, on a
frozen four-object confirmatory split (four independent object batches,
500 candidate labels each, no confirmatory labels used for training or
model selection), the head discriminates well above chance: pooled AUROC
$=0.8896$ ($n=2{,}000$), strongest on Tomato Soup Can ($0.9656$) and
weakest on drill ($0.7854$). Used live, as a candidate-selection method in
its own right, on a fresh 120-scene evaluation disjoint from the head's
own training and dev data, it beats geometry pooled (McNemar's exact
$p<0.0001$, driven mainly by mustard and Tomato Soup Can), but is
statistically indistinguishable from the already-reported neighborhood
mean/point aggregators (pooled $p=0.857$; per object, only mustard reaches
significance against geometry among the head-vs-geometry comparisons, and
none of the head-vs-neighborhood-method comparisons reach significance in
either direction). A related uncertainty-penalized score,
$\mathrm{mean}-0.5{\times}\mathrm{std}$ over the same five perturbations,
shows the identical pattern: it beats geometry (McNemar's exact
$p=0.00012$) but is not distinguishable from the plain neighborhood mean or
worst-case aggregators already in \cref{tab:fourobj_neighborhood} (all
$p\geq1.0$ against those two). We report both as consolidating negatives,
following \cref{sec:critic}'s pairwise-loss-contribution finding: a
learned robustness signal and an uncertainty-penalized heuristic both
recover the same benefit as the simpler neighborhood aggregators already
established above, not a further, separable improvement over them
(\cref{tab:fourobj_negative}).

\begin{table}[t]
\centering
\caption{Consolidating negative results, robustness-head and uncertainty-penalized extensions (\cref{sec:fourobj}), reported with equal rigor to the positive findings above.}
\label{tab:fourobj_negative}
\small
\begin{tabular}{@{}p{0.32\linewidth}p{0.6\linewidth}@{}}
\toprule
Item & Result \\
\midrule
Robustness head vs.\ geometry & Beats geometry pooled ($p{<}0.0001$); driven by mustard ($p{=}0.0051$) and Tomato Soup Can ($p{=}0.0001$) \\
Robustness head vs.\ neighborhood mean/point & Statistically indistinguishable (pooled $p{=}0.857$); no per-object comparison reaches significance either direction \\
Uncertainty-penalized score ($\mathrm{mean}-0.5\mathrm{std}$) vs.\ geometry & Beats geometry ($p{=}0.00012$) \\
Uncertainty-penalized score vs.\ neighborhood mean/worst-case & Statistically indistinguishable ($p\geq1.0$ both) \\
\bottomrule
\end{tabular}
\end{table}

Taken together, this section supports a bounded version of \cref{sec:critic}'s
claim: an object-centric, causally-admissible critic improves candidate
selection and local-perturbation robustness on a fourth object added after
the design was frozen, the benefit is object-dependent (near-ceiling
objects show no measurable gain), it does not generalize zero-shot to an
unseen object class, and neither a learned auxiliary robustness head nor an
uncertainty-penalized heuristic measurably improves on the simplest
neighborhood-aggregation baseline already established. All results in this
section are simulation-only, live-executed unless stated otherwise as
offline re-scoring, consistent with \cref{sec:critic,sec:multihead}'s own
distinction between the two.

% =============================================================================
\section{Case Study: A Silent Gripper-Controller Scaling Bug}
\label{sec:gripperbug}

While investigating why grasp success generalized poorly across object
categories on the Piper platform, we traced the failure one level below
the usual diagnostic stopping point -- not ``where does the arm end up,''
but ``what does the gripper actually command at the actuator level'' --
and found that RoboSuite's composite gripper controller was silently
re-scaling an already-real-units gripper action a second time. The
platform-specific gripper's action-formatting function computes and
returns an absolute, real-units joint-position target directly (a
porting artifact from a different upstream convention), but the composite
controller's default action-scaling behavior applies its own
normalized-to-real-units rescaling on top of it, unconditionally. At a
verified, fully converged ``fully open'' command, the actual commanded
actuator range spanned roughly one tenth of a millimeter, against a
documented $76$--$100\,\mathrm{mm}$. A first fix attempt -- supplying
explicit input-range bounds -- had no effect, traced to the controller's
config-assembly code discarding those keys during rebuild; the actual fix
was a single boolean flag, disabling action scaling for this sub-config,
letting the already-correct action pass through unmodified.

A paired $n{=}20$-per-object pilot validated the fix: Cracker
$0\%\rightarrow45\%$, Mustard $0\%\rightarrow65\%$, Pear unaffected at
$65\%$ (consistent with its prior $70$--$90\%$ range within pilot-scale
noise). A final confirmatory run at paper scale ($n{=}40$--$60$ per
object, on a clean, previously unused trial range) reproduced Cracker
($50\%$, $n{=}40$) and Mustard ($70\%$, $n{=}40$) cleanly, but Pear did
not replicate at its original figure -- three independent 20-trial
batches gave $40\%$, $40\%$, and $50\%$ (pooled $43.3\%$, $n{=}60$),
consistently below the original single-batch $65\%$ (Fisher's exact
$p{=}0.12$). No code drift was found that would explain a genuine
regression; we report the larger, replicated figure ($43.3\%$) as the
confirmed rate rather than the flattering first number, consistent with
this paper's stated discipline.

% =============================================================================
\section{Toward Real-Hardware Deployment}
\label{sec:hardware}

Following our project's scope decision to include real-hardware
capability rather than remain simulation-only, we built an interface
skeleton against the documented vendor SDK for our target 6-DoF arm,
written while the physical arm was not connected to the development
machine. Every hardware-touching method raises an explicit error with a
verification marker rather than guessing an unconfirmed unit or scale
convention, following the same relative-delta-clamped,
replay-not-live-IK design pattern used by this project's validated
SO-ARM101 hardware track. We present this as an architecture and an
itemized pre-deployment safety checklist, not as a hardware result --
physical validation is explicitly future work, gated on the checklist
items stated in the design document.

% =============================================================================
\section{What Doesn't Work: An Honest Negative-Results Ledger}
\label{sec:negative}

\Cref{tab:ledger} summarizes eleven mechanisms tested and closed over the
course of this investigation, each evaluated with a paired-trial design
and McNemar's exact test (or the noted equivalent). Full mechanistic
explanations for each row are given in \cref{app:ledger}.

\begin{table}[t]
\centering
\caption{Ruled-out mechanisms (see \cref{app:ledger} for full detail).}
\label{tab:ledger}
\small
\begin{tabular}{@{}p{0.03\linewidth}p{0.29\linewidth}p{0.6\linewidth}@{}}
\toprule
\# & Mechanism & Result \\
\midrule
1 & OT-coupled flow matching for candidates & $-10.0$pp vs.\ baseline ($p{=}0.0025$); negative on all 7 objects \\
2 & Condition-aware OT fix (C\textsuperscript{2}OT) & Partial, inconsistent mitigation \\
3 & Real-time world-model action correction & Net negative, 3 independent pilots ($-9.3$, $-18.7$, $-13.3$pp) \\
4 & Training-free consensus candidate selection & Effective on SO-ARM101; did not replicate cross-embodiment (re-verified\textsuperscript{*}, null confirmed) \\
5 & AR-mediated active-calibration sampling & Active worse than passive ($p{=}0.0048$); refinement null \\
6 & Cross-embodiment reranking (causally invalid) & $p{<}0.0001{\times}5 \rightarrow$ exact null ($p{=}1.0000$) once corrected -- \cref{sec:method} \\
7 & Binary contact-triggered hard stop & Re-verified\textsuperscript{*}: null ($65\%$ vs.\ $65\%$, $p{=}1.0000$) \\
8 & Sustained-contact stop (5-step) & Re-verified\textsuperscript{*}: null both batches ($p{=}1.0000$) \\
9 & Force-proportional admittance damping & Re-verified\textsuperscript{*}: null, 0 discordant pairs \\
10 & Gripper pre-narrowing before descent & Re-verified\textsuperscript{*}: worse, $75\%$ vs.\ $15\%$, $\mathbf{p{=}0.0005}$ \\
11 & World-model sim-to-real, scarce real data & Not piloted -- closed at feasibility stage \\
\bottomrule
\end{tabular}
\end{table}
\vspace{-0.5em}
{\footnotesize\textsuperscript{*}Row 4 and rows 7--10 were re-run after
\cref{sec:rngbug}'s placement-RNG fix; see below -- the consensus-selection
null (row 4) was reconfirmed under corrected pairing, and three of the four
rows 7--10 negative findings did not replicate.}

Rows 7--10 were originally reported as four independently negative
mechanisms. Re-verifying them under \cref{sec:rngbug}'s placement-RNG fix
told a materially different, more precise story: \textbf{three of the
four (hard stop, sustained-contact stop, admittance damping) have no
measurable effect under valid pairing} -- their original "made things
worse" findings were very likely themselves artifacts of the same RNG bug,
which added spurious between-condition variance that looked like a real
effect (row 9's re-verification is especially stark: $75\%$ vs.\ $75\%$
with \emph{zero} discordant pairs across $n{=}20$, meaning the binary
outcome matched on literally every trial). \textbf{Only row 10
(pre-narrowing gripper clearance before descent) has a genuine, replicated
effect}, and it strengthened substantially under correction ($p{=}0.0005$
vs.\ the original $p{=}0.11$, not significant). The corrected reading:
contact-\emph{response} interventions (react once contact during descent
is detected) are inert for this failure mode, most likely because they
rarely trigger in a way that matters; the one intervention that changes
baseline geometry \emph{before} any contact occurs is the one with a real
cost. This is narrower and more mechanistically precise than "everything
we tried backfired," and it is the reading we stand behind. We state
plainly that the underlying failure mode is diagnosed, not solved, and
requires either physical tactile sensing or a mechanism fundamentally
different from the four tested here.

% =============================================================================
\section{Conclusion}
\label{sec:conclusion}

We set out to improve cross-embodiment grasp reliability and found,
instead, that our own strongest apparent result in that direction was an
artifact of a silent, well-known-in-the-abstract but rarely formalized
class of evaluation error. Rather than treat that as a dead end, we
formalized the constraint it violated, built and validated both a manual
and an automated tool to check for it, and found -- twice -- that even
careful, motivated verification still misses real instances of the
problem it is looking for, until the process is made mechanical.
Alongside this, we report a genuine, validated fix to a previously
undiagnosed execution bug that unlocked most of a robot platform's latent
grasp reliability, an honest account of eleven mechanisms that did not
work for the problem that remains, and an architecture for the
real-hardware phase this project has not yet reached. We believe the
discipline demonstrated across all of these -- verify against ground
truth, report the number that survives replication rather than the one
that looked best first, and say plainly what remains unsolved -- is a
more durable contribution than any single mechanism we tested.

As a next step, motivated by independent evidence \cite{consensusuncertainty2025}
that synthetic Gaussian perturbation does not reproduce real
pose-estimator failure modes, we are instrumenting a real RGB-D sensor to
empirically characterize pose-estimation noise structure (per-axis
variance, Gaussianity, position-orientation correlation) against the
isotropic-Gaussian assumption our own candidate-selection experiments
relied on, and -- contingent on that measurement actually showing a
structural difference, not assumed in advance -- to test whether a
selection rule informed by the measured structure changes the
cross-embodiment generalization result in \cref{tab:ledger} row 4. Any
such feature remains checkable against \cref{sec:method}'s criterion,
since it would be computed purely from the candidate pool's own poses and
the measured noise covariance. This work has not yet started data
collection and is not included in the results reported above.

% =============================================================================
\section*{Acknowledgment}
% DRAFT -- author must verify/expand this against the full writing history of
% the manuscript before submission; this reflects only what was directly
% observed in the sessions that produced the object-relative counterfactual
% critic, mechanistic multi-head analysis, and four-object extension content,
% not necessarily the complete AI-assistance history of the whole paper.
The authors used Claude (Anthropic) to assist with drafting and editing
portions of the manuscript text -- including the object-relative
counterfactual critic, mechanistic multi-head analysis, and four-object
extension sections, and the abstract -- with statistical analysis scripting,
and with literature search and citation verification. All experimental
design, code implementation, data collection, and scientific claims were
conducted and verified by the authors.

% =============================================================================
\appendix
\section{Full Ruled-Out-Mechanisms Ledger}
\label{app:ledger}

Every entry was tested with a paired-trial design and McNemar's exact
test (or the noted equivalent).

\begin{table*}[t]
\centering
\caption{Full mechanistic detail for every ruled-out method.}
\small
\begin{tabular}{@{}p{0.02\linewidth}p{0.19\linewidth}p{0.22\linewidth}p{0.27\linewidth}p{0.24\linewidth}@{}}
\toprule
\# & Method & Mechanism & Result & Why it failed \\
\midrule
1 & OT-CFM candidate generation & Minibatch-OT-coupled conditional flow matching, replacing random CoM sampling & Pooled $-10.0$pp vs.\ baseline, $p{=}0.0025$; negative on all 7 objects & Minibatch OT coupling ignores conditioning (object identity) -- training/inference distribution mismatch \\
2 & C\textsuperscript{2}OT condition-aware OT fix & Condition-aware cost-matrix reweighting & Partial, inconsistent mitigation; hardest object showed no improvement & Fix validated only on generation-quality metrics, not physical execution \\
3 & MPC-style world-model correction & Small-range action search before gripper close & Net negative, 3 pilots ($-9.3$/$-18.7$/$-13.3$pp) & Real-time in-loop correction search consistently made things worse \\
4 & Consensus (median-pose) selection & Pick pool-median candidate, not lowest IK error & Effective on SO-ARM101 (Fisher's exact significant); no cross-embodiment replication under two noise models. Re-verified under \cref{sec:rngbug}'s fix (Pear $80\%/70\%$, Mustard $80\%/70\%$, Cracker $50\%/50\%$, all null) & Training-free heuristic; genuine generalization failure, not a bug, and not an RNG-bug artifact \\
5 & AR-mediated active calibration & Rollout prediction shown via AR; active human-judgment sampling & Active worse than passive ($p{=}0.0048$); refinement null ($p{=}0.32$--$0.73$) & No benefit from active solicitation over passive logging \\
6 & Cross-embodiment reranking & Pooled-data reranker for live candidate selection & $p{<}0.0001{\times}5 \rightarrow$ null after 1st correction; 2nd leak found by auto-tagger $\rightarrow$ null confirmed, accuracy $0.8236{\rightarrow}0.1327$ & No causally valid pre-execution signal predicts outcome; best valid feature: $0.06$ correlation \\
7 & Binary contact hard-stop & Stop descend on any detected contact & \emph{Original} (RNG-bug-affected): worse, $40\%$ vs.\ $70\%$. \textbf{Re-verified} (\cref{sec:rngbug}'s fix): null, $65\%$ vs.\ $65\%$, $p{=}1.0000$, 2 discordant pairs & Original finding does not replicate under valid pairing -- no measurable effect \\
8 & Sustained-contact stop & Require 5 consecutive contact-positive steps & \emph{Original}: promising batch 1 ($75\%/55\%$), opposite-direction batch 2 ($65\%/50\%$), combined $p{=}1.0000$. \textbf{Re-verified}: both batches agree, $75\%$ vs.\ $70\%$, $p{=}1.0000$ each & Original instability was itself likely an RNG-bug artifact; valid pairing gives a stable, consistent null \\
9 & Force-admittance damping & Progress damped by contact force magnitude & \emph{Original}: worse in 2 calibrations ($55\%/75\%$, $50\%/75\%$). \textbf{Re-verified}: $75\%$ vs.\ $75\%$, \emph{zero} discordant pairs across $n{=}20$ & Original mechanistic story (exposure-window extension) does not hold under valid pairing -- no measurable effect, not a harmful one \\
10 & Gripper pre-narrowing & Narrow gripper before descend & \emph{Original}: worse, $20\%$ vs.\ $50\%$, $p{=}0.11$ (n.s.). \textbf{Re-verified}: worse, $15\%$ vs.\ $75\%$, $\mathbf{p{=}0.0005}$ & The only one of the four with a genuine, replicated, strengthened effect -- permanently reducing pre-descent clearance measurably hurts \\
11 & World-model sim-to-real transfer & Train in sim, deploy with little/no real data & Not piloted -- closed at feasibility stage & Working prior art needs infrastructure (thousands of real demos, or video-diffusion-scale backbone) this project lacks \\
\bottomrule
\end{tabular}
\end{table*}

% =============================================================================
\begin{thebibliography}{10}

\bibitem{mahler2017dexnet}
J.~Mahler, J.~Liang, S.~Niyaz, M.~Laskey, R.~Doan, X.~Liu, J.~A.~Ojea,
and K.~Goldberg,
``Dex-Net 2.0: Deep Learning to Plan Robust Grasps with Synthetic
  Point Clouds and Analytic Grasp Metrics,''
\emph{Proc.\ Robotics: Sci.\ Syst.\ (RSS)}, 2017.

\bibitem{graspgen2025}
A.~Murali, B.~Sundaralingam, Y.-W.~Chao, W.~Yuan, J.~Yamada,
M.~Carlson, F.~Ramos, S.~Birchfield, D.~Fox, and C.~Eppner,
``GraspGen: A Diffusion-based Framework for 6-DOF Grasping with
  On-Generator Training,''
\emph{arXiv:2507.13097}, 2025.

\bibitem{graspgenx2026}
B.~Han, Y.-W.~Chao, E.~Coumans, C.~Eppner, B.~Sundaralingam,
J.~Deng, S.~Birchfield, and A.~Murali,
``GraspGen-X: Cross-Embodiment 6-DOF Diffusion-based Grasping,''
\emph{arXiv:2606.00998}, 2026.

\bibitem{fpl2026}
M.~Torne, A.~Mahajan, A.~Bhat, and C.~Finn,
``Freeform Preference Learning for Robotic Manipulation,''
\emph{arXiv:2606.32027}, 2026.

\bibitem{feasibility2025}
T.~Liu, R.~Dazeley, B.~Champion, and A.~Cosgun,
``Optimizing Robotic Placement via Grasp-Dependent Feasibility
  Prediction,''
\emph{arXiv:2512.18922}, 2025.

\bibitem{offlinerl_leakage2026}
H.~Wang, J.~Bowden, C.~Crosby, and S.~Bansal,
``Offline Policy Evaluation for Manipulation Policies via Discounted
  Liveness Formulation,''
\emph{arXiv:2605.11479}, 2026.

\bibitem{coorgrasp2026}
M.~Yu, Y.~Jiang, Y.~Jia, R.~Yi, and X.~Li,
``CoorGrasp: Coordinated Contact Control for Adaptive Dexterous
  Grasping Under Uncertainty,''
\emph{arXiv:2607.03557}, 2026.

\bibitem{rothgrammar2026}
S.~Roth,
``A Grammar of Machine Learning Workflows: Rejecting Data Leakage at
  Call Time,''
\emph{arXiv:2603.10742}, 2026.

\bibitem{leakagedetector2025}
E.~A.~AlOmar, C.~DeMario, R.~Shagawat, and B.~Kreiser,
``LeakageDetector: An Open Source Data Leakage Analysis Tool in
  Machine Learning Pipelines,''
\emph{arXiv:2503.14723}, 2025.

\bibitem{benchmarkaudit2026}
T.~Jiang, X.~Tan, S.~Wheeler, L.~Sun, T.~W.~Ayalew, and M.~Walter,
``What Are We Actually Benchmarking in Robot Manipulation?,''
\emph{arXiv:2606.04233}, 2026.

\bibitem{consensusuncertainty2025}
E.~C.~Joyce, Q.~Zhao, N.~Burgdorfer, L.~Wang, and P.~Mordohai,
``Consensus-Driven Uncertainty for Robotic Grasping based on RGB
  Perception,''
\emph{Proc.\ IEEE/RSJ Int.\ Conf.\ Intelligent Robots and Systems (IROS)},
2025.

\bibitem{lggsn_ral_under_review}
%% Self-citation placeholder -- this project's own SO-ARM101/LGGSN paper,
%% submitted to RA-L 2026-07-11, under review at time of writing. Do NOT
%% cite as an accepted/published reference until the review outcome is
%% known; update venue/year once resolved.
Anonymous,
``When Geometry Is Not Enough: Object-Dependent Failure Modes of
  Learning-to-Rank Grasp Selection in Open-World Manipulation,''
\emph{Under review}, 2026.

\end{thebibliography}

\end{document}
